\subsection{Checking the session}
\label{sec:pfit-check}

After creating a session, the user runs

\begin{verbatim}
/pfit-check <session-name>
\end{verbatim}

This step verifies that the problem is ready for translation and fitting. It
reads the generated configuration, the generated model file, and the CSV files
named in the configuration. Its purpose is to catch errors in the scientific
specification before numerical optimization begins.

The check answers four practical questions. First, are the data files usable by
the framework: rectangular CSV files, valid time columns, expected column names,
and compatible missing-value conventions? Second, is the model specification
self-consistent: declared states, parameters, initial conditions, observables,
and data columns all refer to the same objects? Third, is the loss well defined
for the available measurements, including blank or \texttt{NaN} entries? Fourth,
are the requested numerical settings adequate for the scale and structure of
the problem?

The agent may revise the generated YAML or Python files when the correction is
a mechanical consequence of the specification. It must not modify the data
files. Any change to a CSV file is the user's responsibility, because such a
change alters the measurements being fitted.

The report is written to

\begin{verbatim}
sessions/<session-name>/generated/user_input_check.txt
\end{verbatim}

and separates findings into critical errors, warnings, and recommendations.
Critical errors are expected to prevent translation or a valid run and must be
resolved before continuing. Warnings identify choices that may make the fit less
reliable or harder to interpret. Recommendations are generated from the
framework's documented checking rules and must be supported by evidence from
the current session, such as parameter ranges, data scales, missing values,
solver settings, or loss construction. Exwaitamples include using a masked loss for
missing measurements, increasing the global search budget, tightening solver
tolerances for small state variables, or normalizing residuals across measured
channels.

Scientific ambiguities are reported to the user rather than guessed. For
example, the checker may ask whether a blank measured value is intentional,
whether an unused data column should be ignored, or whether a mismatch between a
configured initial condition and the first data row reflects the model or the
file. After resolving such issues, the user reruns \texttt{/pfit-check}. When no
critical errors remain, the session is ready for JAX translation:

\begin{verbatim}
/pfit-jax <session-name>
\end{verbatim}
